% ICSE NIER: four pages of main text, fifth page references only.
% Evidence: structural_study/DISCOVERY_REPORT.md and results/discovery_validation.json.
% No figure, paid generation, or new subject execution accompanies this revision.
\RequirePackage[T1]{fontenc}
\documentclass[10pt,conference]{IEEEtran}
\usepackage{amsmath,amssymb}
\usepackage{booktabs,array,tabularx}
\usepackage{xurl}
\usepackage[hidelinks]{hyperref}
\newcommand{\code}[1]{{\urlstyle{tt}\nolinkurl{#1}}}
\newcommand{\tracer}{\textsc{Tracer}}
\title{Ghost Scripts in Agent Trajectories}
\author{}
\begin{document}
\maketitle

\begin{abstract}
Coding agents can leave behind scripts that constructed a repair but are no longer needed by the evaluated task. We study these ghost scripts through static analysis of recorded trajectories. Abstract syntax trees (ASTs) identify source-writing behavior, while invocations and code versions establish construction relationships. Parsing embedded source reveals whether retained implementations match the final product or have been superseded. We introduce \tracer{} (Trace-based Runtime Analysis for Code Evolution Review), an agent-independent architecture linking recorded reasoning to symbol changes, revisions, and a decision ledger. Across 88 passing mini-SWE-agent runs on SWE-bench, the analysis finds 65 source-writing candidates among 1,662 Python file occurrences. History corroborates 62 helpers, recovering an existing positive reference set whose saved deletion checks establish task-level redundancy. Of 25 helpers with embedded definitions linked to product code, 20 retain current implementations and five retain superseded ones. Efficiency and security analysis are future applications of the broader architecture.
\end{abstract}
\begin{IEEEkeywords}
coding agents, static analysis, software provenance, code evolution
\end{IEEEkeywords}

\section{Introduction}
After a coding agent repairs a method, the script that rewrote it may remain. Its string literals can contain an implementation later replaced in the product. A reviewer needs to establish their relationship and the script's remaining purpose. Passing task tests alone does not answer those questions.

Agents such as SWE-agent and OpenHands edit and execute code~\cite{sweagent,openhands,codeact}, producing programs that construct other programs. A lack of repository callers cannot distinguish a discarded script from a maintained shell-invoked utility. Ordinary ASTs also treat embedded implementations as string data. Understanding these artifacts requires their development context.

Our idea is to analyze trajectories as histories of code construction. Static analysis identifies source writes; recorded invocations and workspace versions associate them with product changes. Parsing embedded source and following corresponding definitions then establishes implementation currency. We call a retained construction helper a \emph{ghost script} when separate deletion evidence shows that it is redundant under a fixed task oracle.

Reconstructing such context after development becomes a software archaeology task. \tracer{}, \emph{Trace-based Runtime Analysis for Code Evolution Review}, records it as development occurs. This agent-independent architecture links prompts and recorded reasoning to individual file and symbol changes through Git diffs, preserves revisions, and derives a decision ledger. The tool supports analyses of how code came to be; the mini-SWE-agent integration and AST study on SWE-bench~\cite{swebench} are one application.

We present \tracer{}'s architecture and a static analysis of its recorded history for identifying construction helpers and characterizing their code. Across 88 passing runs, the analysis recovers 62 reference helpers with saved deletion evidence. Five of 25 helpers linked to product definitions retain superseded implementations. These results establish retrospective feasibility for this application; broader tool benefits remain to be evaluated.

\section{What Survives an Agent Session?}
\label{sec:concept}
In the recorded repair of \code{django__django-13346} by MiniMax M2.5, the agent invokes \code{/tmp/patch2.py} to write a \code{KeyTransformIn} class into \code{django/db/models/fields/json.py}. Later actions replace it. At the passing endpoint, the product uses \code{batch_process_rhs}, while the retained script's \code{new_class} string contains the earlier \code{resolve_expression_parameter} and \code{process_rhs} methods.

The helper's ordinary AST exposes writing operations but no executable class or function definitions. Parsing its string data reveals the class. Linking it to the invocation and subsequent product versions distinguishes a retained construction attempt from the endpoint implementation.

We distinguish three properties. \emph{Construction role} associates a surviving helper with a source-writing action. \emph{Currency} describes whether a linked implementation matches the endpoint, has been superseded, or cannot be resolved. \emph{Task-level redundancy} requires deletion evidence. For final workspace $F$, helper $h$, and the outcome vector $v_O$ of fixed oracle $O$, the latter is
\begin{equation}
v_O(F\setminus\{h\})=v_O(F).
\label{eq:redundancy}
\end{equation}
Both evaluations must be valid, contain the required checks, and use the same oracle. The equality does not establish that $h$ has no future maintenance value. It also does not follow from AST similarity, a missing caller, or absence of later execution.

\section{TRACER Architecture}
\label{sec:tracer}
\tracer{} separates agent adapters, code-linked recording, and downstream analyses. Its unit is a development step with an identifier, prompt, recorded reasoning, raw events, author, and before/after version references. The adapter contract is independent of the agent or model. Existing interfaces accept OpenHands events and explicit human edits; another agent requires an adapter to the same step format.

\subsection{Linking Reasoning to Symbol Changes}
The structural index comes from the separate \code{codebase-memory-mcp} project~\cite{codebasememory}. \tracer{} reads its code graph, including qualified symbol names, source spans, and relationships. For each Git revision range, changed line ranges are mapped to the innermost overlapping symbols. Live AST extraction resolves new symbols and changed spans; re-indexing and reconciliation update their graph identities. Unresolved ranges fall back to file-level links.

These links attach the step's prompt and recorded reasoning to its affected code. Provenance tags distinguish narration, tool-call thoughts, exposed chain of thought (CoT), and provider summaries when available. A step touching several symbols links to each; the attachment does not establish a separate causal explanation for every edit. Unexposed reasoning is unavailable, and recorded text is not assumed faithful intent.

\subsection{Revisions and the Decision Ledger}
The Git-backed driver checkpoints development steps, retaining intermediate attempts. A separate SQLite trace store holds steps, node links, signatures, and symbol co-change records. Per-version snapshots preserve structural changes. Development lineages retain fork points and merge visibility, while historical records remain inspectable after later edits. Optional reasoning embeddings support search; they are not required for recording.

The decision ledger is derived from recorded revisions without an LLM. For each step, it compares added line text with session-end removals through two Git diffs. Normalized-text overlap yields \emph{survived}, \emph{revised}, or \emph{abandoned} entries. Symbol entries additionally check whether linked definitions remain. Error observations followed by further actions generate \emph{error\_retry} entries. Records retain abandoned content and heuristic links to later surviving steps on the same symbol or file. These are reproducible development summaries, not correctness verdicts or ghost labels.

\subsection{Reusable History and Analysis Interfaces}
File/symbol history queries, version navigation, and review bundles expose diffs with their recorded reasoning and ledger entries. The panel connects code structure to the step timeline. The intended release separates this general \tracer{} tool from agent connectors and study-specific analyses. Determinism concerns deriving records and reconstructing captured versions from stored evidence, not repeating model outputs or guaranteeing complete runtime capture.

\section{AST Analysis of Retained Helpers}
\label{sec:method}
This application has two stages, identifying construction helpers and characterizing their retained code. Both consume recorded evidence offline; deletion outcomes are joined afterward.

\subsection{Connecting mini-SWE-agent}
The study adapter records commands, exit codes, raw and model-visible output, and model requests/responses. Git-anchored, content-addressed snapshots preserve initial and post-action workspaces, including untracked files and bounded temporary-file captures. It exports the \tracer{} step shape using snapshot identifiers instead of per-action Git commits. The AST analysis consumes these records directly. It evaluates neither CoT explanations nor the general decision ledger; helper links and definition currency are derived by the analysis below.

\subsection{Recovering Source-Writing Behavior}
For each retained Python path, we analyze every captured version. Abstract values track constant paths, file handles, read origins, and transformations. Assignments propagate values through \code{open}, \code{pathlib} operations, string/regular-expression replacements, and selected list mutations. A source-writing candidate requires a resolved Python destination and records the written value's origins; filenames and complete embedded functions are not required.

This bounded possible-effect analysis joins branches, approximates loops once, and binds arguments for supported direct local calls. Deferred bodies are not assumed executed. Dynamic paths and unsupported flows remain unresolved. The analysis covers common helper idioms without a soundness guarantee for arbitrary Python.

\subsection{Corroborating the Construction History}
Shell parsing locates direct Python-script invocations, separates commands from heredoc bodies and quoted arguments, and resolves supported working-directory changes. The invoked version must have identical before/after bytes or match a quoted body written immediately before invocation.

The static write destination must match a repository Python file changed between the action's snapshots. A \emph{history-supported construction relationship} retains the helper version, read origin, write site, target versions, and event identifier. Compound actions establish interval association rather than isolated causation.

The target need not have two parseable, different ASTs. A helper may change comments, introduce a syntax error, or repair an invalid intermediate file. We record target parse status and AST fingerprints separately from byte changes. Helper revisions receive their own AST transitions so that a role observed for an earlier version is not silently attributed to its final bytes.

\subsection{Following Embedded Implementations}
After discovery, a separate characterization stage extracts string literals containing source definitions. An artificial statement wrapper permits indented fragments without changing whitespace inside multiline string values. We fingerprint functions, classes, and asynchronous functions using location-free AST serialization, preserving names, literals, operators, decorators, and docstrings.

Let $H$ be the AST fingerprint and $D_e(f,q)$ the definition at file $f$ and qualified name $q$ in state $e$. An embedded payload $p$ is a produced-version candidate when
\begin{equation}
H(p)=H(D_{e^+}(f,q)),\qquad
H(p)\ne H(D_{e^-}(f,q)),
\label{eq:link}
\end{equation}
where $e^-$ and $e^+$ denote the action's before/after states. The payload matches a definition observed after the action but absent or different before it. Known absence has a distinguished fingerprint; unavailable source is not absence. Whole-class matches subsume contained methods. Argument-use tags and source inspection distinguish replacement payloads from search patterns and unused literals.

We follow definitions satisfying Equation~\ref{eq:link} through later snapshots. Equal endpoint fingerprints indicate current correspondence; a unique different definition indicates supersession. Missing or renamed definitions, ambiguous duplicates, and parse failures remain unresolved. This preserves reversions and avoids treating every edit to the enclosing file as a change to the linked definition. The output is an evidence record, with deletion outcomes attached separately where available.

\section{Exploratory Study}
\label{sec:study}
We ask \emph{RQ1: Can static analysis of agent trajectories identify retained construction helpers? RQ2: What is the nature of their retained code, and which helpers have evidence of task-level redundancy?}

\paragraph{Population}
The archive contains 240 runs using mini-SWE-agent 2.4.6~\cite{mini}, with 12 model configurations on 20 SWE-bench Verified tasks across eight repositories. The models are DeepSeek V3-0324, V3.1-Terminus, V3.2; MiniMax M2, M2.1, M2.5; Qwen3-Coder-Next, Qwen3.5-35B-A3B, Qwen3.5-122B-A10B; and Llama 3.1-70B-Instruct, 3.3-70B-Instruct, 4-Maverick. All used temperature 0.2 and the official agent templates; provider reasoning defaults were not controlled.

The cohort comprises 88 endpoints with matching passing check vectors across two saved evaluations. Ten configurations contribute; Llama 3.1 and Maverick contribute none. The exclusions are 149 recorded nonpassing runs, one invalid endpoint, and two unavailable replay environments. Exact model IDs and per-model counts accompany the artifact.

The full captured inventory contains 1,695 $(\text{run},\text{path})$ occurrences across 56 eligible runs, including 1,411 pytest-temporary artifacts. The other 32 runs contribute no entry. Python analysis covers 1,662 files; 33 are unsupported. Eleven Python endpoints fail parsing, all test-generated.

\paragraph{Reference and validation}
Discovery discards prior nomination, origin, and deletion labels. We subsequently compare with 62 frozen positive references and join saved deletion results by run, path, and hash. These examples informed development, including invocation refinements after reference misses. The comparison measures retrospective recovery, not population precision or recall.

No new agent runs or subject executions were performed. Hash and snapshot checks plus 26 local tests cover analysis boundaries and a synthetic archive without nomination inputs. Source inspections were Codex-assisted; independent human adjudication remains future work.

\subsection{RQ1: Identification and Historical Evidence}
AST analysis finds 65 source-writing candidates. History corroborates 62 files across 40 runs: 27 repository files and 35 external temporary files (Table~\ref{tab:discovery}). These recover every reference helper and its nominated edit event. The AST-only stage already includes all 62 references; history adds evidence about actual construction intervals rather than increasing reference recall.

\begin{table}[t]
\centering
\caption{Discovery over the full captured inventory. Counts are file occurrences; the reference comparison is retrospective.}
\label{tab:discovery}
\begin{tabular}{@{}lr@{}}
\toprule
Stage & Files \\
\midrule
Retained source inventory & 1,695 \\
Python files & 1,662 \\
AST source-writing candidates & 65 \\
History-supported helpers & 62 \\
Positive-reference helpers recovered & 62/62 \\
Additional AST candidates, deletion unknown & 3 \\
\bottomrule
\end{tabular}
\end{table}

The 62 helpers participate in 67 edit intervals. Literal creation establishes the invoked version in 54 intervals; 13 have equal before/after helper blobs. Every tracked read origin matches the write target. Three intervals involve older helper versions. Six helpers have multiple versions, contributing seven AST changes.

Target transitions include 56 different parseable AST pairs, one byte change with an equal AST, and ten with an unparseable intermediate target. Construction roles remain observable when outputs are later repaired.

The three additional candidates have unknown deletion status. Two, in Xarray 3677 and pytest 5262, are invoked but leave identical target blobs. The third, in Sphinx 8056, is a retained copy whose original is executed. They are not labeled non-ghosts.

The initial prototype corroborated 58 references. Resolving four helpers rewritten and invoked within one action required matching quoted creation bodies to archived bytes. This refinement used the study data.

\subsection{RQ2: Nature of Retained Scripts}
Characterization uses the recovered reference edit events and previously audited comparisons. Twenty-seven helpers contain parseable embedded definitions. Matching links 25 to product definitions across 20 runs and eight tasks. Of 37 unlinked helpers, 35 lack complete parseable definitions, one contains a partial replacement, and one an unused proposed class. Writing analysis still detects them.

Three of 28 raw matches are excluded as pre-existing search patterns or unchanged methods. The 25 primary links use 18 string replacements, four regular-expression replacements, two list insertions, and one list-slice replacement. None of these helpers has executable function/class definitions; their implementations exist in string data.

\begin{table}[t]
\centering
\caption{Currency of the 62 history-supported helpers. All have separate saved evidence of task-oracle redundancy.}
\label{tab:currency}
\begin{tabular}{@{}lrrrr@{}}
\toprule
Location & Helpers & Current & Superseded & Unlinked \\
\midrule
Repository & 27 & 7 & 2 & 18 \\
External temporary & 35 & 13 & 3 & 19 \\
\midrule
Total & 62 & 20 & 5 & 37 \\
\bottomrule
\end{tabular}
\end{table}

Twenty linked helpers match the endpoint; five retain superseded versions across five runs and three tasks (Table~\ref{tab:currency}). These precede changes to Django's \code{KeyTransformIn}, pytest's \code{EncodedFile.mode}, and SymPy tuple-printing logic. Differences remain after stripping docstrings, but do not establish incorrect behavior.

Two Django/Sphinx histories end current after supersession and reversion. Another Django edit changes the enclosing file without changing the class. Endpoint comparison and entity granularity matter.

Saved individual deletion trials establish Equation~\ref{eq:redundancy} for all 62 helpers. Combined deletions preserve baseline vectors in 40 runs across two checks. The audit verifies evaluation validity, matching oracle identifiers, and complete required checks. The redundant set spans all currency categories. The three extra candidates lack deletion evidence; an indispensable-helper reference group is unavailable.

\section{Related Work}
Agent platforms expose editing and execution~\cite{sweagent,openhands,codeact}. RepoCoder and Repoformer retrieve repository context~\cite{repocoder,repoformer}, a possible consumer of construction annotations whose effectiveness we have not evaluated. Git AI records agent authorship through checkpoints and Git metadata~\cite{gitai}; here we examine retained artifacts' roles and relationships to produced code.

ChangeDistiller and GumTree extract changes through tree differencing~\cite{changedistiller,gumtree,gumtree2024}. Clone genealogy follows related fragments across versions~\cite{genealogy}, including in human--agent development~\cite{aigenealogy}. Our AST fingerprints connect source stored as a helper's data to a product definition and follow it through the repair history.

Provenance research distinguishes entities, activities, and derivations~\cite{prov}. noWorkflow combines AST analysis and runtime information~\cite{noworkflow}, YesWorkflow exposes script workflows through annotations~\cite{yesworkflow}, and ReproZip captures execution dependencies~\cite{reprozip}. Our analysis uses existing action snapshots to interpret retained construction artifacts without instrumenting or rerunning them.

Delta debugging isolates relevant input or changes~\cite{dd}, Perses uses syntax to guide reduction~\cite{perses}, and Chisel addresses program debloating~\cite{chisel}. TRIM minimizes agent-generated patches using trajectories and reports removal of scratch scripts~\cite{trim}. Here, construction provenance describes a helper's writing role and represented version. Separate deletion checks validate redundancy; automatic removal remains future work.

\section{Threats to Validity}
The study is retrospective and conditions on passing, reconstructable runs from one scaffold and a small task sample. Tasks and models recur, so files are not independent replications. Recovering 62 known positives does not establish general detection accuracy, especially because development used those cases. Unselected files have not been independently labeled as non-helpers.

The inventory covers captured new files, not every repository symbol; temporary capture is incomplete. Results concern workspaces, including external temporary files, rather than necessarily submitted patches. Bounded AST and shell analyses miss indirect invocation, reflection, complex APIs, and dynamic paths. Snapshot intervals can contain multiple operations or conceal changes undone within an action. Exact embedded matching misses transformed fragments, while named-definition tracking leaves moves and duplicates unresolved. Assistant inspection and unit tests establish inspectability and internal consistency, not independent empirical validation. Finally, the same task oracle defines eligibility and redundancy; passing deletion checks does not establish general removal safety or future usefulness.

\section{Future Plans}
We will release \tracer{} with separate adapter and analysis packages, then evaluate a frozen helper detector on independently labeled sessions. Maintained generators should be included when comparing precision, recall, and abstention against AST-only screening.

\paragraph{Development efficiency}
Repeated AST states and invocations with unchanged targets nominate work for inspection. Relating these intervals to recorded duration and request usage could quantify the cost of repeated attempts. Independent labels must distinguish avoidable work from useful exploration and verification before testing efficiency interventions.

\paragraph{Security analysis}
History could support review of security-relevant changes, such as the removal and restoration of input validation, by linking edits to surrounding commands and versions. Independently adjudicated cases should test reviewers' ability to locate introducing and corrective edits against a final diff and an unlinked trajectory. Vulnerability and exploitability require separate evidence.

\paragraph{Validated cleanup in mini-SWE-agent}
An end-of-session stage could nominate helpers, archive their history, and remove them under a retention policy with validation. Evaluation must include indispensable helpers and measure incorrect removal, task success, and overhead. These interventions remain unevaluated.

\section{Conclusion}
\tracer{} preserves development history for later analyses. This AST application recovers 62 reference helpers with saved redundancy evidence and distinguishes current from superseded implementations. Efficiency and security studies are future applications of the broader architecture.

\section*{Artifact Availability}
The release will separate \tracer{} core, the mini-SWE-agent adapter, and AST analyses, with model configurations and reproduction evidence. Packaging and public deposition remain pending.

\clearpage
\begin{thebibliography}{22}
\bibitem{sweagent}
J.~Yang, C.~E. Jimenez, A.~Wettig, K.~Lieret, S.~Yao, K.~Narasimhan, and O.~Press, \textquotedblleft SWE-agent: Agent-computer interfaces enable automated software engineering,\textquotedblright\ in \emph{NeurIPS}, 2024. [Online]. Available: \url{https://arxiv.org/abs/2405.15793}
\bibitem{openhands}
X.~Wang \emph{et al.}, \textquotedblleft OpenHands: An open platform for AI software developers as generalist agents,\textquotedblright\ in \emph{ICLR}, 2025. [Online]. Available: \url{https://openreview.net/forum?id=OJd3ayDDoF}
\bibitem{codeact}
X.~Wang, Y.~Chen, L.~Yuan, Y.~Zhang, Y.~Li, H.~Peng, and H.~Ji, \textquotedblleft Executable code actions elicit better LLM agents,\textquotedblright\ in \emph{ICML}, 2024. [Online]. Available: \url{https://proceedings.mlr.press/v235/wang24h.html}
\bibitem{swebench}
C.~E. Jimenez, J.~Yang, A.~Wettig, S.~Yao, K.~Pei, O.~Press, and K.~Narasimhan, \textquotedblleft SWE-bench: Can language models resolve real-world GitHub issues?\textquotedblright\ in \emph{ICLR}, 2024. [Online]. Available: \url{https://arxiv.org/abs/2310.06770}
\bibitem{codebasememory}
DeusData, \textquotedblleft codebase-memory-mcp,\textquotedblright\ software repository. [Online]. Available: \url{https://github.com/DeusData/codebase-memory-mcp}. Accessed: Oct. 9, 2026.
\bibitem{mini}
The SWE-agent Team, \textquotedblleft mini-SWE-agent documentation.\textquotedblright\ [Online]. Available: \url{https://mini-swe-agent.com/latest/}. Accessed: Oct. 9, 2026.
\bibitem{repocoder}
F.~Zhang \emph{et al.}, \textquotedblleft RepoCoder: Repository-level code completion through iterative retrieval and generation,\textquotedblright\ in \emph{EMNLP}, 2023. [Online]. Available: \url{https://aclanthology.org/2023.emnlp-main.151/}
\bibitem{repoformer}
D.~Wu, W.~U. Ahmad, D.~Zhang, M.~K. Ramanathan, and X.~Ma, \textquotedblleft Repoformer: Selective retrieval for repository-level code completion,\textquotedblright\ in \emph{ICML}, 2024. [Online]. Available: \url{https://proceedings.mlr.press/v235/wu24a.html}
\bibitem{gitai}
Git AI, \textquotedblleft Agent blame: See which AI agent wrote every line of code.\textquotedblright\ [Online]. Available: \url{https://usegitai.com/agent-blame}. Accessed: Oct. 9, 2026.
\bibitem{changedistiller}
B.~Fluri, M.~W\"ursch, M.~Pinzger, and H.~C. Gall, \textquotedblleft Change distilling: Tree differencing for fine-grained source code change extraction,\textquotedblright\ \emph{IEEE Trans. Softw. Eng.}, vol.~33, no.~11, pp. 725--743, 2007.
\bibitem{gumtree}
J.-R.~Falleri, F.~Morandat, X.~Blanc, M.~Martinez, and M.~Monperrus, \textquotedblleft Fine-grained and accurate source code differencing,\textquotedblright\ in \emph{ASE}, 2014, pp. 313--324, doi: 10.1145/2642937.2642982.
\bibitem{gumtree2024}
J.-R.~Falleri and M.~Martinez, \textquotedblleft Fine-grained, accurate and scalable source differencing,\textquotedblright\ in \emph{ICSE}, 2024, Art. 231, doi: 10.1145/3597503.3639148.
\bibitem{genealogy}
M.~Kim, V.~Sazawal, D.~Notkin, and G.~C. Murphy, \textquotedblleft An empirical study of code clone genealogies,\textquotedblright\ in \emph{ESEC/FSE}, 2005, pp. 187--196, doi: 10.1145/1081706.1081737.
\bibitem{aigenealogy}
D.~Sousa, I.~Uchoa, M.~Paixao, C.~Ragkhitwetsagul, and T.~L. Matos, \textquotedblleft An empirical study of code clone genealogies in human-AI collaborative development,\textquotedblright\ in \emph{MSR}, 2026, pp. 979--983, doi: 10.1145/3793302.3793613.
\bibitem{prov}
Y.~Gil and S.~Miles, Eds., \textquotedblleft PROV model primer,\textquotedblright\ W3C Working Group Note, 2013. [Online]. Available: \url{https://www.w3.org/TR/prov-primer/}
\bibitem{noworkflow}
L.~Murta, V.~Braganholo, F.~Chirigati, D.~Koop, and J.~Freire, \textquotedblleft noWorkflow: Capturing and analyzing provenance of scripts,\textquotedblright\ in \emph{IPAW 2014}, LNCS, vol.~8628. Springer, 2015, pp. 71--83. [Online]. Available: \url{https://braganholo.github.io/papers/murta2014-ipaw.pdf}
\bibitem{yesworkflow}
T.~McPhillips \emph{et al.}, \textquotedblleft YesWorkflow: A user-oriented, language-independent tool for recovering workflow information from scripts,\textquotedblright\ \emph{Int. J. Digital Curation}, vol.~10, no.~1, pp. 298--313, 2015, doi: 10.2218/ijdc.v10i1.370.
\bibitem{reprozip}
F.~Chirigati, D.~Shasha, and J.~Freire, \textquotedblleft ReproZip: Using provenance to support computational reproducibility,\textquotedblright\ in \emph{TaPP}, 2013. [Online]. Available: \url{https://www.usenix.org/conference/tapp13/technical-sessions/presentation/chirigati}
\bibitem{dd}
A.~Zeller and R.~Hildebrandt, \textquotedblleft Simplifying and isolating failure-inducing input,\textquotedblright\ \emph{IEEE Trans. Softw. Eng.}, vol.~28, no.~2, pp. 183--200, 2002.
\bibitem{perses}
C.~Sun, Y.~Li, Q.~Zhang, T.~Gu, and Z.~Su, \textquotedblleft Perses: Syntax-guided program reduction,\textquotedblright\ in \emph{ICSE}, 2018, pp. 361--371, doi: 10.1145/3180155.3180236.
\bibitem{chisel}
K.~Heo, W.~Lee, P.~Pashakhanloo, and M.~Naik, \textquotedblleft Effective program debloating via reinforcement learning,\textquotedblright\ in \emph{ACM CCS}, 2018, pp. 380--394, doi: 10.1145/3243734.3243838.
\bibitem{trim}
A.~Mathai, S.~Iyer, A.~Nogikh, P.~Maniatis, F.~Ivancic, J.~Yang, and B.~Ray, \textquotedblleft TRIM: Reducing AI-generated CodeSlop via agent trajectory minimization,\textquotedblright\ arXiv preprint arXiv:2607.18161, 2026.
\end{thebibliography}
\end{document}
